\documentclass[../main.tex]{subfiles}
\begin{document}
\chapter{Introduction}
\section{Motivational Problem}
Suppose that we wish to find the shortest path between 2 points:
\begin{enumerate}
  \item in Euclidean space?
  \item on the surface of a sphere?
  \item on any general curved surface, say in $\R^{3}$?
\end{enumerate}
For \textbf{(i)}, we could write the path as $y = y(x)$ for $a \leq x \leq b$ and compute its length using the \textit{functional} $L[y]$:
\[
  L[y] = \int_{a}^{b} \sqrt{1 + (y')^2} \d{x}
\]
and can produce similar functionals for computing the length in the other cases.
A \textit{functional} is a map which takes a function $y(x)$ to a real number.
Since the above function computes the length of the curve $y(x)$, our original problem reduces to wanting know which function $y(x)$ minimises this functional.

In this course, we will develop tools to allow us to determine these minima and maxima of functionals.
\section{\texorpdfstring{Classifying Stationary Points in $\R^{n}$}{Classifying Stationary Points in n-dimensions}} \label{classifyingStatPts}
Consider a twice continuously function $f: \R^{n} \to \R$ and general point as $\vec{x} = (x_1, \ldots, x_n) \in \R^{n}$.
Assuming that $f$ is at least twice continuously differentiable, we say that $\vec{a} \in \R^{n}$ is a \textit{stationary point} of $f$ if $\nabla f(\vec{a}) = \vec{0}$.

We can then investigate the behaviour around such points using a Taylor expansion:
\[
  f(\vec{x}) = f(\vec{a}) + \sum_{i = 1}^{n} (x_i - a_i)\cancelto{0}{(\nabla_i f(\vec{a}))} + \frac{1}{2}\sum_{i, j = 1}^{n} (x_i - a_i)(x_j - a_j)H_{ij}(\vec{a}) + O(|\vec{x} - \vec{a}|^3)
\]
where $H_{i j}(\vec{a}) = \frac{\partial^2 f}{\partial x_i \partial x_j}(\vec{a})$.
\begin{remark}[Recap]
  Recall from IA Differential Equations that $H_{i j}$ is the \textit{Hessian matrix} for $f$.
  This is symmetric (i.e. $H_{i j} = H_{j i}$) since partial derivatives commute.
\end{remark}

Without loss in generality, we can assume that $\vec{a} = \vec{0}$ (as otherwise we can just translate) and so we have:
\[
  f(\vec{x}) - f(\vec{0}) = \frac{1}{2}x_i H_{i j}(\vec{0})x_j + O(|\vec{x}|^3) \quad  \text{(using $\Sigma$ convention)}
\]
Since $H_{i j}(\vec{0})$ is symmetric, we know from IA Vectors and Matrices that we can carry out an orthogonal diagonalisation.
Further, we can pick our diagonalisation matrix so that it is a rotation as we can simply reorder its columns (and thus also the columns the diagonalised form) to ensure that $\det R = 1$.
The components of vectors are then transformed by $x_i = R_{i j}x_j'$ and $H'(\vec{0}) = R^{\trans} H(\vec{0}) R$ is the diagonal matrix:
\[
  H_{i j}'(\vec{0}) = \begin{pmatrix}
  \lambda_1 &  &  \\
   & \ddots &  \\
   &  & \lambda_n \\
  \end{pmatrix}
\]
where $\lambda_i$ are the eigenvalues of $H$, which, since $H$ is real and symmetric, must be real.

Therefore in this new basis we have:
\[
  f(\vec{x}) - f(\vec{0}) = \frac{1}{2}\lambda_i (x_i')^2 + O(|\vec{x}'|^3)
\]
This allows us to use the eigenvalues to determine the nature of the stationary point at $\vec{x} = \vec{0}$:
\begin{itemize}
  \item If $\lambda_i > 0\ \forall i \in \{1, \ldots, n\}$, then $f(\vec{x}) > f(\vec{0})$ for small enough $|\vec{x}|$, i.e. $\vec{x} = \vec{0}$ is a local minimum as the function increases in all directions around $\vec{x} = \vec{0}$.
  \item Similarly, if $\lambda_i < 0\ \forall i \in \{1, \ldots, n\}$, then $f(\vec{x}) < f(\vec{0})$ for small enough $|\vec{x}|$, i.e. $\vec{x} = \vec{0}$ is a local maximum as the functions decreases in all directions around $\vec{x} = \vec{0}$.
  \item If the eigenvalues are non-zero but have mixed signs, then we have a saddle point as the function increases in one direction but decreases in another for sufficiently small $|\vec{x}|$.
  \item If any of the eigenvalues are zero then it is a \textit{degenerate stationary point}.
    The behaviour at these points then depends on behaviour depends on $O(|\vec{x}'|^3)$ terms.
\end{itemize}
\begin{remark}[In 2D]
  For $n = 2$, we know that $\tr H = \lambda_1 + \lambda_2$ and $\det H = \lambda_1 \lambda_2$.
  Thus:
  \begin{itemize}
    \item If $\det H > 0$ and $\tr H > 0$, then we have a local minimum.
    \item If $\det H > 0$ and $\tr H < 0$, then we have a local maximum.
    \item If $\det H < 0$, then we have a saddle point as the eigenvalues are non-zero and have opposite signs.
    \item If $\det H = 0$, then one of the eigenvalues is $0$ and so this is a degenerate case.
  \end{itemize}
\end{remark}
\begin{remark}
  If we have a \textit{local} minimum that is \textbf{not} a \textit{global} minimum, then either:
  \begin{itemize}
    \item There is another stationary point that is the global minimum.
    \item The function is defined only on a subset of $\R^{n}$ and so the global minimum could occur on the boundary of the domain which would not require it to be a stationary point.
    \item The function has no global minimum, for example, it is unbounded below.
  \end{itemize}
  We can make analogous remarks in the case of maxima.
\end{remark}
\begin{example}
  Consider the function $f: \R^2 \to \R,\ (x, y) \mapsto x^3 + y^3 - 3xy$.
  We see that:
  \begin{align*}
    \nabla f = \begin{pmatrix}
    3x^2 - 3y \\
    3y^2 - 3x \\
  \end{pmatrix} = \vec{0} &\iff x^2 = y \text{ and } y^2 = x \\
  &\implies y^4 = y \\
  &\implies y = 0 \text{ and } x = 1 \\
  &\ \ \text{ or } y = 1 \text{ and } x = 0 \text{ since $y \in \R$}
  \end{align*}
  So the stationary points are at $(0, 0)$ and $(1, 1)$.

  We can also find the hessian $H$:
  \[
    H = \begin{pmatrix}
    6x & -3 \\
    -3 & 6y \\
    \end{pmatrix} \implies \det H = 9(4xy - 1) \text{ and } \tr H = 6(x + y)
  \]
  Thus we can conclude that:
  \begin{itemize}
    \item At $(1, 1)$, $\det H = 27 > 0$ and $\tr H = 12 > 0$, so it is a local minimum with $f = - 1$.
    \item At $(0, 0)$, $\det H = -9 < 0$, so it is a saddle point with $f = 0$.
  \end{itemize}
  The directions for the saddle point can be found using the eigenvectors of the hessian:
  \begin{align*}
    &\lambda_1 = - 3 \text{ with } \vec{e}_1 = \begin{pmatrix}
    1 \\
    1 \\
    \end{pmatrix} \\
    &\lambda_2 = 3 \text{ with } \vec{e}_2 = \begin{pmatrix}
    1 \\
    -1 \\
    \end{pmatrix}
  \end{align*}
  Thus $(0, 0)$ is a local maximum along $y = x$ but is a local minimum along $y = -x$.

  Note also that for large negative or positive values of $x$ and $y$, the function decreases or increases without bound and so it does not have a global minimum or maximum.
\end{example}
\section{Convex Functions}
\subsection{Defining Convexity}
\begin{definition}[Convex Set]
  A set $S \subseteq \R^{n}$ is called \textit{convex} if for all $\vec{x}, \vec{y} \in S$ and for every $t \in (0, 1)$, the point $t \vec{x} + (1 - t)\vec{y}$ is also in $S$.
\end{definition}
\begin{remark}[Intuition]
  A set is convex if for any two points, every point on the line segment joining these to points is also contained within the set.
\end{remark}
\begin{definition}[Graph]
  Let $f$ be a function with domain $D(f) \subseteq \R^{n}$, i.e. $f: D(f) \to \R$.
  The \textit{graph} of $f$ is the surface $z = f(\vec{x})$ in $\R^{n + 1}$ with coordinates $(\vec{x}, z)$.
\end{definition}
\begin{definition}[Chord]
  A \textit{chord} of $f$ is a line segment joining two points on a graph.
\end{definition}
\begin{definition}[Convex Function]
  A function $f: D(f) \to \R$ is \textit{convex} if both:
  \begin{enumerate}
    \item The set $D(f)$ is convex.
    \item The graph of $f$ lies below, or on, any of its chords.
      More precisely, for any $\vec{x}, \vec{y} \in D(f)$ and all $t \in (0, 1)$:
      \begin{equation}
        f(t \vec{x} + (1 - t)\vec{y}) \leq t f(\vec{x}) + (1 - t) f(\vec{y}) \label{convexFn}
      \end{equation}
      where \textbf{(i)} ensures that $t\vec{x} + (1 - t)\vec{y}$ is always in the domain of $f$.
  \end{enumerate}
\end{definition}
\begin{definition}[Strictly Convex Function]
  A function $f$ is \textit{strictly convex} if it is convex but we instead adapt \textbf{(ii)} above so that the inequality in \cref{convexFn} is strict:
  \begin{equation}
    f(t \vec{x} + (1 - t)\vec{y}) < t f(\vec{x}) + (1 - t) f(\vec{y})\ \forall \vec{x} \neq \vec{y} \label{strictConvexFn}
  \end{equation}
  i.e. the graph of $f$ always lies below and \textbf{not on} its chords.
\end{definition}
\begin{remark}
  We say that $f$ is \textit{concave} if $-f$ is convex and similarly for the strict case.
\end{remark}
\begin{example}
  \begin{enumerate}
    \item Consider $f(x) = x^2$.
      We have that:
      \[
        [(1 - t)x + ty]^2 - (1 - t)x^2 - ty^2 = - t(1 - t)(x - y) < 0
      \]
      since $t \in (0, 1)$ and $x \neq y$.
      Thus $f$ is strictly convex.
    \item Consider $f(x) = e^{x}$ and $g(x) = e^{-x}$, these are both strictly convex functions.
    \item Consider $f(x) = |x|$.
      This is convex as:
      \begin{align*}
        |(1 - t)x + ty| - (1 - t)|x| - t|y| &\leq |1 - t||x| + |t||y| - (1 - t)|x| - t|y| \\
                                            &= 0
      \end{align*}
      However, it is not strictly convex as if we pick two points on the same branch, then their chord lies on the graph.
  \end{enumerate}
\end{example}
\subsection{First Order Conditions}
\begin{theorem}
  If $f$ is differentiable then convexity is equivalent to:
  \[
    f(\vec{y}) \geq f(\vec{x}) + (\vec{y} - \vec{x}) \cdot \nabla f(\vec{x})\ \forall \vec{x}, \vec{y} \in D(f)
  \]
  That is, the graph of $f$ lies on or above any of its tangent planes.
\end{theorem}
\begin{proof}
  \begin{proofdirection}{Assume $f$ is convex}
    For any fixed $\vec{x}, \vec{y} \in D(f)$ define:
    \[
      h(t) = (1 - t)f(\vec{x}) + tf(\vec{y}) - f((1 - t)\vec{x} + t\vec{y})
    \]
    From IA Vector Calculus we know that:
    \[
      \deriv{}{t}[f((1 - t)\vec{x} + t\vec{y})] = (\vec{y} - \vec{x}) \cdot \nabla f((1 - t)\vec{x} + t \vec{y})
    \]
    Thus we have:
    \[
      h'(0) = - f(\vec{x}) + f(\vec{y}) - (\vec{y} - \vec{x}) \cdot \nabla f(\vec{x})
    \]
    and so our claim is equivalent to showing that $h'(0) \geq 0$.
    \[
      h'(0) = \lim_{t \to 0} \frac{h(t) - h(0)}{t} = \lim_{t \to 0} \frac{h(t)}{t}
    \]
    Since $f$ is convex, $h(t) \geq 0\ \forall t \in (0, 1) \implies h'(0) \geq 0$ and so the result follows.
  \end{proofdirection}
  \begin{proofdirection}{Assume the inequality holds}
    For any $\vec{x}, \vec{y}, \vec{z} \in D(f)$ we have:
    \begin{align*}
      f(\vec{x}) &\geq f(\vec{z}) + (\vec{x} - \vec{z}) \cdot \nabla f(\vec{z}) \\
      f(\vec{y}) &\geq f(\vec{z}) + (\vec{y} - \vec{z}) \cdot \nabla f(\vec{z})
    \end{align*}
    For $t \in (0, 1)$, multiplying the first inequality by $(1 - t)$ and the second by $t$ and summing we have:
    \[
      (1 - t)f(\vec{x}) + tf(\vec{y}) \geq f(\vec{z}) + [(1 - t)\vec{x} + t\vec{y} -  \vec{z}] \cdot \nabla f(\vec{z})
    \]
    Since $\vec{z}$ is arbitrary, we can pick $\vec{z} = (1 - t)\vec{x} + t\vec{y}$ to yield:
    \[
      (1 - t)f(\vec{x}) + tf(\vec{y}) \geq f((1 - t)\vec{x} + t\vec{y})\ \forall \vec{x}, \vec{y} \in D(f), \ t \in (0, 1)
    \]
    and so $f$ is convex.
  \end{proofdirection}
\end{proof}
\begin{corollary}
  If $f$ is convex and differentiable, then any stationary point is a global minimum. \label{globalMin}
\end{corollary}
\begin{proof}
  If $\vec{a}$ is a stationary point, then $\nabla f(\vec{\vec{a}}) = \vec{0}$ so:
  \[
    f(\vec{y}) \geq f(\vec{\vec{a}})\ \forall \vec{y} \in D(f)
  \]
  i.e. there is a global minimum at $\vec{a}$.
\end{proof}
\begin{proposition}[Alternative First Order Conditon]
  An equivalent first order condition is: \label{altFirstOrderCond}
  \begin{equation}
    (\vec{y} - \vec{x}) \cdot (\nabla f(\vec{y}) - \nabla f(\vec{x})) \geq 0\ \forall \vec{x}, \vec{y} \in D(f) \label{equiv1stEqn}
  \end{equation}
\end{proposition}
\begin{remark}
  For $n = 1$ this becomes $(y - x)(f'(y) - f'(x)) \geq 0$ and so $y \geq x \implies f'(y) \geq f'(x)$, i.e. $f'$ is monotonically increasing.
\end{remark}
\begin{proof}
  \begin{proofdirection}{Assume the original first order condition holds}
    We have:
    \begin{align*}
      f(\vec{y}) &\geq f(\vec{x}) + (\vec{y} - \vec{x}) \cdot \nabla f(\vec{x}) \iff -(\vec{y} - \vec{x}) \cdot \nabla f(\vec{x}) \geq f(\vec{x}) - f(\vec{y}) \\
      f(\vec{x}) &\geq f(\vec{y}) + (\vec{x} - \vec{y}) \cdot \nabla f(\vec{y}) \iff (\vec{y} - \vec{x}) \cdot \nabla f(\vec{y}) \geq f(\vec{y}) - f(\vec{x})
    \end{align*}
    Taking the sum of these two inequalities, the result follows.
  \end{proofdirection}
  \begin{proofdirection}{Assume assume that the new first order condition holds}
    Let $\vec{z}(t) = (1 - t)\vec{x} + t\vec{y}$ and consider:
    \begin{align}
      f(\vec{y}) - f(\vec{x}) &= \eval{f(\vec{z}(t))}{t = 0}{t = 1} \nonumber \\
                              &= \int_{0}^{1} \deriv{}{t}[f(\vec{z}(t))] \d{t} \nonumber \\
                              &= \int_{0}^{1} (\vec{y} - \vec{x}) \cdot \nabla f(\vec{z}(t)) \d{t} \nonumber \\
      \implies f(\vec{y}) - f(\vec{x}) - (\vec{y} - \vec{x}) \cdot \nabla f(\vec{x}) &= \int_{0}^{1} (\vec{y} - \vec{x}) \cdot [\nabla f(\vec{z}(t)) - \nabla f(\vec{x})] \d{t} \tag{$\dag$} \label{eqn1stDag}
    \end{align}
    Applying \cref{equiv1stEqn} with $\vec{x}$ and $\vec{z}$ we have:
    \begin{align*}
      (\vec{z}(t) - \vec{x}) \cdot (\nabla f(\vec{z}(t)) - \nabla f(\vec{x})) &\geq 0 \\
      \iff t(\vec{y} - \vec{x}) \cdot (\nabla f(\vec{z}(t)) - \nabla f(\vec{x})) &\geq 0 \\
      \iff \vec{y} - \vec{x} \cdot (\nabla f(\vec{z}(t)) - \nabla f(\vec{x})) &\geq 0 \text{ as $t \in (0, 1)$}
    \end{align*}
    Thus the integrand on the right hand side of \cref{eqn1stDag} is always non-negative and so:
    \[
      f(\vec{y}) \geq f(\vec{x}) + (\vec{y} - \vec{x}) \cdot \nabla f(\vec{x})
    \]
  \end{proofdirection}
\end{proof}
\section{Second Order Conditions}
\begin{theorem}[Second Order Condition for Convexity]
  If $f$ has a continuous second derivative then convexity is equivalent to all eigenvalues of the Hessian $H(\vec{x})$ being non-negative for all $\vec{x} \in D(f)$, i.e. $H(\vec{x})$ is positive semi-definite for all $\vec{x} \in D(f)$. \label{2ndOrderConvexity}
\end{theorem}
\begin{remark}[Recap]
  Recall that a symmetric matrix $A$ is positive semi-definite if $\vec{x}^{\trans} A \vec{x} \geq 0\ \forall \vec{x} \in \R^{n}$. This is also equivalent to all of its eigenvalues being non-negative.
\end{remark}
\begin{proof}
  \begin{proofdirection}{Assume $f$ is convex}
    Using \cref{altFirstOrderCond}, setting $\vec{y} = \vec{x} + \vec{h}$:
    \[
      \vec{h} \cdot [\nabla f(\vec{x} + \vec{h}) - \nabla f(\vec{x})] \geq 0
    \]
    Using a Taylor expansion:
    \[
      \nabla_i f(\vec{x} + \vec{h}) = \nabla_i f(\vec{x}) + h_j H_{i j}(\vec{x}) + O(|\vec{h}|^3)
    \]
    and so:
    \[
      h_i H_{i j}(\vec{x}) h_j + O(|\vec{h}|^3) \geq 0
    \]
    Suppose for contradiction that $H(\vec{x})$ has a negative eigenvalue $\lambda$ with eigenvector $\vec{e}$ then we could set $\vec{h} = h \vec{e}$ giving:
    \[
      \lambda h^2 |\vec{e}|^2 + O(|\vec{h}|^3) \geq 0
    \]
    but this should hold for arbitrary $h$ and $\lambda < 0$ so for small enough $h$ the left hand side is negative and so we have a contradiction.
  \end{proofdirection}
  \begin{proofdirection}{Assume the second order condition holds}
    For ease, in this course we will only consider the case of $n = 1$.
    In this case $H(x) = f''(x)$ so the second order condition tells us $f''(x) \geq 0\ \forall x \in D(f)$.
    We then have:
    \begin{align*}
      0 &\leq \operatorname{sign}(y - x) \int_{x}^{y} f''(z) \d{z} \\
        &= \operatorname{sign}(y - x) (f'(y) - f'(x))
    \end{align*}
    since the sign of the integral and $(y - x)$ always disagree.
    Multiplying both sides by $|y - x|$ we have:
    \[
      0 \leq (y - x)(f'(y) - f'(x))
    \]
    since $|y - x| \operatorname{sign}(y - x) = y - x$.
    This is \cref{altFirstOrderCond} in the case of $n = 1$ and so we are done.
  \end{proofdirection}
\end{proof}
\begin{example}
  Consider the function:
  \[
    f(x, y) = \frac{1}{xy} \text{ with } D(f) = \{(x, y) \in \R^2: x > 0, y > 0\}
  \]
  We first note that $D(f)$ is a convex set, else there would be no chance that $f$ could be convex.
  Computing the Hessian we have:
  \[
    H = \frac{1}{xy} \begin{pmatrix}
    \frac{2}{x^2} & \frac{1}{xy} \\
    \frac{1}{xy} & \frac{1}{y2} \\
    \end{pmatrix}
  \]
  We see that $\det H = \frac{3}{x^{4}y^4} > 0$ and $\tr H = \frac{2}{xy}(\frac{1}{x^2} + \frac{1}{y^2}) > 0$.
  Thus the eigenvalues of $H$ are strictly positive everywhere in $D(f)$ and so the second order condition tells us that it is convex.

  We might be tempted to instead set:
  \[
    D(f) = \{(x, y) \in \R^{2}: xy > 0\}
  \]
  as the Hessian would still be positive semi-definite everywhere within this this domain.
  However this domain is not convex and so $f$ cannot be convex, despite the Hessian conditions still holding.
\end{example}
\end{document}
